% !TEX root = ../Projektdokumentation.tex
\section{Entwurfsphase} 
\label{sec:Entwurfsphase}

\subsection{Zielplattform}
\label{sec:Zielplattform}
Im Rahmen der Entwurfsphase wurde die Zielplattform für die Umsetzung des Systems festgelegt. Dabei wurden sowohl die
Anforderungen des Projekts als auch die im schulischen Umfeld verfügbaren Technologien berücksichtigt. Die Anwendung wird als Desktop-Anwendung
entwickelt und basiert auf der Programmiersprache Python. Python wurde gewählt, da im schulischen Unterricht größtenteils diese Sprache verwendet
wurde und diese somit eine sichere und effiziente Umsetzung ermöglicht. Zudem bietet Python eine große Anzahl an Bibliotheken, die die Entwicklung unterstützen.

Für die grafische Benutzeroberfläche wird die Bibliothek Tkinter eingesetzt. Tkinter ist Bestandteil der Standardbibliothek von
Python und ermöglicht die Erstellung einfacher, plattformunabhängiger Benutzeroberflächen. Aufgrund der guten Integration in Python
und der geringen Einstiegshürde eignet sich Tkinter insbesondere für den Einsatz im schulischen Kontext.
Zur Datenhaltung wird ein Microsoft SQL Server verwendet. Dieser stellt ein relationales Datenbankmanagementsystem dar und
ermöglicht eine strukturierte und konsistente Speicherung der benötigten Daten, wie beispielsweise Benutzerinformationen, Bestellungen
und Produktbestände. Die Verwaltung der Datenbank erfolgt über das Tool SQL Server Management Studio 22.
Die Kommunikation zwischen Anwendung und Datenbank wird über die Python-Bibliothek \emph{pyodbc} realisiert. Diese ermöglicht den
Zugriff auf relationale Datenbanken über eine standardisierte Schnittstelle (\acs{ODBC}) und unterstützt die Ausführung von \acs{SQL}-Abfragen direkt
aus der Anwendung heraus.

Durch die gewählte Kombination aus Python, Tkinter und Microsoft SQL Server wird eine klare Trennung zwischen Benutzeroberfläche, Anwendungslogik
und Datenhaltung erreicht. Gleichzeitig ermöglicht die Zielplattform eine stabile und nachvollziehbare Umsetzung der Anforderungen im Rahmen des Projekts.

\subsection{Architekturdesign}
\label{sec:Architekturdesign}
Für den Aufbau der Anwendung wurde eine an das Model-View-Controller-Muster (\acs{MVC}) angelehnte Softwarearchitektur gewählt. Ziel dieser Architektur
ist die Trennung von Benutzeroberfläche, Anwendungslogik und Datenzugriff, um die Anwendung übersichtlicher, wartbarer und besser erweiterbar zu gestalten.

Das \textbf{Model} umfasst in diesem Projekt insbesondere den Zugriff auf die Datenbank sowie die Verarbeitung und Bereitstellung der benötigten Daten.
Hierzu zählen unter anderem Funktionen zum Abrufen von Bestellungen, Guthaben, Benachrichtigungen und Produktbeständen sowie zum Speichern von Änderungen,
beispielsweise bei Stornierungen oder Guthabenbuchungen.

Die \textbf{View} bildet die grafische Benutzeroberfläche ab. Diese wurde mit der Python-Bibliothek Tkinter umgesetzt und stellt die verschiedenen Benutzerseiten,
Dialoge und Steuerelemente bereit. Dazu gehören unter anderem Tabellen zur Anzeige offener Bestellungen, Eingabefelder für Ticketnummern, Schaltflächen zur Navigation
sowie die Darstellung des \acs{QR}-Codes für die Abholung.

Der \textbf{Controller} übernimmt die Steuerung der Benutzerinteraktionen und verbindet die Benutzeroberfläche mit der Datenlogik. Er verarbeitet Eingaben, ruft
Datenbankfunktionen auf und aktualisiert anschließend die dargestellten Inhalte in der Oberfläche. Dadurch wird sichergestellt, dass Änderungen im System, wie
zum Beispiel neue Bestellungen oder Statusänderungen, direkt in der Anwendung sichtbar werden.

Die Umsetzung orientiert sich an diesem Muster, ohne es vollständig in strikt getrennte Klassenstrukturen zu überführen. Insbesondere innerhalb einzelner
\acs{GUI}-Klassen werden sowohl Darstellungslogik als auch steuernde Funktionen zusammengefasst. Dennoch ist eine fachliche Trennung der Verantwortlichkeiten erkennbar:
Die Oberfläche dient primär der Darstellung, während Datenbankoperationen über separate Methoden und Objekte ausgeführt werden.
Ein Beispiel hierfür ist die Besucheransicht. Diese stellt dem Benutzer Informationen wie Ticketnummer, Guthaben, offene Bestellungen, Benachrichtigungen und
den \acs{QR}-Code grafisch zur Verfügung. Gleichzeitig werden Benutzeraktionen, wie das Stornieren einer Bestellung, das Freischalten einer Bestellung für Freund*innen
oder das Aufladen von Guthaben, über entsprechende Steuerungsfunktionen verarbeitet. Die dafür notwendigen Daten werden über die Datenzugriffsschicht aus der Datenbank
geladen und nach Änderungen erneut in der Oberfläche dargestellt.

Durch die Wahl einer an \acs{MVC} angelehnten Architektur wird die Anwendung logisch strukturiert. Dies erleichtert die Entwicklung im Projektverlauf sowie die spätere
Wartung und Erweiterung der Software. Insbesondere die Trennung von Oberflächenkomponenten und Datenzugriff trägt dazu bei, Änderungen gezielter umsetzen zu können.

\subsection{Entwurf der Benutzeroberfläche}
\label{sec:Benutzeroberflaeche} 
Im Rahmen der Entwurfsphase wurde die Benutzeroberfläche der Anwendung zunächst in Form von Mockups geplant. Ziel dieser Vorab-Visualisierung war es,
die Struktur, Navigation und grundlegende Bedienlogik der Anwendung festzulegen, bevor mit der eigentlichen Implementierung begonnen wurde.
Die Mockups dienen dabei als konzeptionelle Grundlage für die spätere Umsetzung und ermöglichen eine frühzeitige Absprache der Benutzerführung mit dem Auftraggeber
sowie der Anordnung der einzelnen Bedienelemente. Insbesondere wurde darauf geachtet, die Benutzeroberfläche übersichtlich und intuitiv zu gestalten, um eine einfache
Bedienung für alle Benutzergruppen zu gewährleisten.

Für jede zentrale Benutzerrolle (Besucher*innen, Standpersonal) wurden entsprechende Oberflächen entworfen. Dabei wurden typische Nutzungsszenarien berücksichtigt, wie
beispielsweise das Aufgeben einer Bestellung, die Anzeige von Bestellstatus oder die Verwaltung von Bestellungen durch das Standpersonal.
Ein besonderer Fokus lag auf der klaren Darstellung von Statusinformationen, wie beispielsweise dem Fortschritt einer Bestellung. Darüber hinaus wurden wichtige Funktionen
wie das Stornieren von Bestellungen oder das Aufladen von Guthaben leicht zugänglich gestaltet. Die im Entwurfsprozess erstellten Mockups sind im Anhang unter
\bhref{app:Mockups} einsehbar und dienten als Vorlage für die spätere programmatische Umsetzung.
Während der Implementierungsphase wurden kleinere Anpassungen vorgenommen, um technische Gegebenheiten sowie praktische Anforderungen zu berücksichtigen.
Durch den Einsatz von Mockups konnte die Benutzeroberfläche bereits vor der Implementierung strukturiert geplant und potenzielle Usability-Probleme frühzeitig erkannt werden.

\subsection{Datenmodell}
\label{sec:Datenmodell}
Im Rahmen der Entwurfsphase wurde ein Datenmodell entwickelt, um die für das System benötigten Informationen strukturiert und konsistent in einer relationalen Datenbank
abzubilden. Ziel war es, alle relevanten Datenobjekte des Bestell- und Abholprozesses sowie deren Beziehungen zueinander nachvollziehbar zu modellieren. Als Grundlage
für die Modellierung wurde zunächst ein Entity-Relationship-Diagramm erstellt. Dieses dient der konzeptionellen Darstellung der wichtigsten Entitäten, Attribute und Beziehungen.
Aufbauend darauf wurde ein relationales Datenbankmodell entworfen, welches die technische Umsetzung in Microsoft SQL Server beschreibt. Beide Darstellungen sind im Anhang unter
\bhref{app:ER-Diagramm} und \bhref{app:Relationales_Datenbankmodell} einsehbar.

Das Datenmodell bildet die zentralen Bestandteile des Systems ab. Dazu zählen insbesondere die Entitäten \emph{Tickets}, \emph{Orders}, \emph{OrderPositions},
\emph{Products}, \emph{Stands}, \emph{Status} und \emph{Messages}. Ergänzt werden diese durch Verknüpfungstabellen wie \emph{Stand2Product} und
\emph{Order2Ticket}, welche n:m-Beziehungen zwischen den jeweiligen Entitäten auflösen.

Bei der Modellierung wurde darauf geachtet, Redundanzen weitgehend zu vermeiden und Beziehungen zwischen den Entitäten über Primär- und Fremdschlüssel eindeutig abzubilden. Das
Datenbankmodell befindet sich in der dritten Normalform.  Durch die relationale Struktur können Daten konsistent gespeichert und effizient abgefragt werden. Dies bildet die
Grundlage für die zuverlässige Verarbeitung von Bestellungen, Beständen und Benutzerinformationen innerhalb der Anwendung.

\subsection{Geschäftslogik}
\label{sec:Geschaeftslogik}
Um die Umsetzung der Geschäftsobjekte vorab grob planen zu können, wurde während der Entwurfsphase ein Klassendiagramm erstellt. Es macht den ungefähren Aufbau der zu
implementierenden Objektklassen anschaulich und ist im Anhang unter \bhref{app:Klassendiagramm} zu finden.

Mit Blick auf die in Abschnitt \bhref{sec:Architekturdesign} erwähnten Rollen des \acs{MVC}-Musters, nimmt die Klasse \emph{Database} die Rolle des Models ein. Die jeweiligen
Seitenobjekte (\emph{LoginPage}, \emph{VisitorPage}, \emph{OrderPage}, \emph{SellerPage}) bilden die Controller, welche die zugrunde liegende Logik hinter den View-Objekten 
bereitstellen. Über die jeweilige \emph{get\_page()} Methode wird innerhalb der Controller-Klassen das, auf Tkinter aufgebaute, View-Objekt initialisiert und mit dem Controller
verbunden.

Für den Dialog zum aufladen des Ticket-gebunden Guthabens (\emph{CreditDialog}) wird eine eigene, wiederverwendbare Klasse implementiert. So kann der Dialog einfach und erweiterbar
an diversen Stellen der Anwendung eingebunden werden.

\subsection{Maßnahmen zur Qualitätssicherung}
\label{sec:Qualitaetssicherung}
Es werden Black-Box-Tests durch eine Test-Gruppe bestehend aus Mitschülern geplant. Die Testfälle werden aus dem zur Verfügung gestellten Anforderungsdokument erstellt. In Folge der
Durchführung der Tests werden Fehler analysiert und korrigiert. Es werden kontinuierliche Feature-Tests durch die Entwickler während der Entwicklung durchgeführt.

\subsection{Pflichtenheft}
\label{sec:Pflichtenheft}
Am Ende der Entwurfsphase wurde ein Pflichtenheft erstellt. Dieses baut auf dem Lastenheft (Abschnitt \bhref{sec:Lastenheft}) auf und beschreibt wie und womit die Autoren die
Anforderungen des Fachbereiches umsetzen möchten. Es dient als Grundlage für die Realisierung des Projektes. Ein Auszug aus dem Pflichtenheft ist im Anhang \bhref{app:Pflichtenheft}
einzusehen.